\documentclass[10pt,a4paper]{amsart}
\usepackage{amsmath,amssymb,amsthm,graphicx}
% For commutative diagrams we use Xy-pic
\usepackage[all]{xy}

% We use 2cell for 2-commutative diagrams.
\xyoption{2cell}
% Theorem environments.
%
\theoremstyle{plain}
\newtheorem{theorem}[subsection]{Theorem}
\newtheorem{proposition}[subsection]{Proposition}
\newtheorem{lemma}[subsection]{Lemma}

\theoremstyle{definition}
\newtheorem{definition}[subsection]{Definition}
\newtheorem{example}[subsection]{Example}
\newtheorem{exercise}[subsection]{Exercise}
\newtheorem{situation}[subsection]{Situation}

\theoremstyle{remark}
\newtheorem{remark}[subsection]{Remark}
\newtheorem{remarks}[subsection]{Remarks}

\numberwithin{equation}{subsection}

% Macros
%
\def\lim{\mathop{\mathrm{lim}}\nolimits}
\def\colim{\mathop{\mathrm{colim}}\nolimits}
\def\Spec{\mathop{\mathrm{Spec}}}
\def\Hom{\mathop{\mathrm{Hom}}\nolimits}
\def\Ext{\mathop{\mathrm{Ext}}\nolimits}
\def\SheafHom{\mathop{\mathcal{H}\!\mathit{om}}\nolimits}
\def\SheafExt{\mathop{\mathcal{E}\!\mathit{xt}}\nolimits}
\def\Sch{\mathit{Sch}}
\def\Mor{\mathop{\mathrm{Mor}}\nolimits}
\def\Ob{\mathop{\mathrm{Ob}}\nolimits}
\def\Sh{\mathop{\mathit{Sh}}\nolimits}
\def\NL{\mathop{N\!L}\nolimits}
\def\CH{\mathop{\mathrm{CH}}\nolimits}
\def\proetale{{pro\text{-}\acute{e}tale}}
\def\etale{{\acute{e}tale}}
\def\QCoh{\mathit{QCoh}}
\def\Ker{\mathop{\mathrm{Ker}}}
\def\Im{\mathop{\mathrm{Im}}}
\def\Coker{\mathop{\mathrm{Coker}}}
\def\Coim{\mathop{\mathrm{Coim}}}

% Boxtimes
%
\DeclareMathSymbol{\boxtimes}{\mathbin}{AMSa}{"02}

%
% Macros for moduli stacks/spaces
%
\def\QCohstack{\mathcal{QC}\!\mathit{oh}}
\def\Cohstack{\mathcal{C}\!\mathit{oh}}
\def\Spacesstack{\mathcal{S}\!\mathit{paces}}
\def\Quotfunctor{\mathrm{Quot}}
\def\Hilbfunctor{\mathrm{Hilb}}
\def\Curvesstack{\mathcal{C}\!\mathit{urves}}
\def\Polarizedstack{\mathcal{P}\!\mathit{olarized}}
\def\Complexesstack{\mathcal{C}\!\mathit{omplexes}}
% \Pic is the operator that assigns to X its picard group, usage \Pic(X)
% \Picardstack_{X/B} denotes the Picard stack of X over B
% \Picardfunctor_{X/B} denotes the Picard functor of X over B
\def\Pic{\mathop{\mathrm{Pic}}\nolimits}
\def\Picardstack{\mathcal{P}\!\mathit{ic}}
\def\Picardfunctor{\mathrm{Pic}}
\def\Deformationcategory{\mathcal{D}\!\mathit{ef}}
\usepackage{cleveref}
\crefname{originalreference}{original reference}{original references}
\Crefname{originalreference}{Original reference}{Original references}
\makeatletter
\newlabel{lemma-constructible-abelian}{{03RZ}{source}{}{sourcepointer-1}{}}
\newlabel{lemma-constructible-abelian@cref}{{[originalreference][0][]03RZ}{[1][0][]source}{}{}{}}
\makeatother
\crefname{originalreference}{original reference}{original references}
\Crefname{originalreference}{Original reference}{Original references}
\makeatletter
\newlabel{lemma-torsion-colimit-constructible}{{03SA}{source}{}{sourcepointer-2}{}}
\newlabel{lemma-torsion-colimit-constructible@cref}{{[originalreference][0][]03SA}{[1][0][]source}{}{}{}}
\makeatother
\crefname{originallemma}{original lemma}{original lemmas}
\Crefname{originallemma}{Original lemma}{Original lemmas}
\makeatletter
\newlabel{properties-lemma-quasi-compact-quasi-separated-spectral}{{094L}{source}{}{sourcepointer-3}{}}
\newlabel{properties-lemma-quasi-compact-quasi-separated-spectral@cref}{{[originallemma][0][]094L}{[1][0][]source}{}{}{}}
\makeatother
\crefname{originallemma}{original lemma}{original lemmas}
\Crefname{originallemma}{Original lemma}{Original lemmas}
\makeatletter
\newlabel{sites-modules-lemma-module-filtered-colimit-constructibles}{{093C}{source}{}{sourcepointer-4}{}}
\newlabel{sites-modules-lemma-module-filtered-colimit-constructibles@cref}{{[originallemma][0][]093C}{[1][0][]source}{}{}{}}
\makeatother
\crefname{originallemma}{original lemma}{original lemmas}
\Crefname{originallemma}{Original lemma}{Original lemmas}
\makeatletter
\newlabel{topology-lemma-constructible-hausdorff-quasi-compact}{{0901}{source}{}{sourcepointer-5}{}}
\newlabel{topology-lemma-constructible-hausdorff-quasi-compact@cref}{{[originallemma][0][]0901}{[1][0][]source}{}{}{}}
\makeatother
\title{Constructible sheaves of Lambda-modules are exactly the cokernels of maps between finite sums of extensions by zero of the constant sheaf}
\author{The Stacks Project Authors}
\date{stacks-project:etale-cohomology.tex:local-snapshot-20261011; GNU FDL 1.2 or later}
\begin{document}
\maketitle
\noindent\textit{Editorial note.} All delivered mathematical text is the author's own, copied from the pinned Stacks Project chapter file etale-cohomology.tex; no statement or proof step was rewritten, reordered or removed.\par\medskip
\noindent\textit{Editorial note.} The counted proof is the author's complete argument: the proof is written out in full: the inclusion from the torsion-colimit lemma, the reduction to a surjection onto a given constructible module, the construction of that surjection from the filtered-colimit description, and the two concluding refinements (cokernel form and affineness of the objects).\par\medskip
\noindent\textit{Editorial note.} The label lemma-category-constructible-modules names the counted statement, and the proof environment immediately following it in the source is delivered verbatim as its proof.\par\medskip
\noindent\textit{Editorial note.} The counted statement is carried by the author's own lemma environment in the chapter source.\par\medskip
\noindent\textit{Editorial note.} The statement and proof are copied from the pinned chapter file; the five cited inputs are kept as external pointers with their canonical Stacks tags.\par\medskip
\noindent\textit{Editorial note.} The preamble is the chapter's own native preamble excerpt already used for this source: the Xy-pic setup (native preamble lines 14-18) and the native theorem declarations with the mathematical macros (native preamble lines 154-217).\par\medskip
\noindent\textit{Allowed standard prerequisites (editorial).}\begin{itemize}
\item Constructible sheaves of Lambda-modules on X\_etale. Same chapter, the definition of constructibility for sheaves of Lambda-modules on the etale site of a quasi-compact quasi-separated scheme. Supplies the objects on both sides of the asserted identification.
\item Abelianness of the category of constructible modules and the colimit description. Same chapter, the lemmas that constructible modules form an abelian category and that every constructible module is a filtered colimit of constructible submodules, together with Modules on Sites, the representation of such a module as a surjective image of a finite sum of extensions by zero. Supplies the surjection from a finite sum of j\_\{U\_i!\}Lambda\_\{U\_i\} onto a given constructible module F and the abelian structure used to conclude.
\item Affineness of the members of an etale covering. Properties, the spectrality of the underlying space of a quasi-compact quasi-separated scheme and the topology lemma allowing the objects to be taken affine. Supplies the final claim that the U\_i and V\_j may be assumed affine.
\end{itemize}\medskip
\ifdefined\hypertarget\hypertarget{sourcepointer-1}{}\fi
\noindent\textit{Source reference.} Original reference 03RZ: sources/stacks-etale-cohomology/pinned-original-source.tex, source lines 11828--11828. Prerequisite pointer invoked at proof line(s) 12431; it names a standard fact of the same chapter that the counted argument uses, and every step of the counted proof built on it is retained verbatim in the delivered body.\par\medskip
\ifdefined\hypertarget\hypertarget{sourcepointer-2}{}\fi
\noindent\textit{Source reference.} Original reference 03SA: sources/stacks-etale-cohomology/pinned-original-source.tex, source lines 12216--12216. Prerequisite pointer invoked at proof line(s) 12428; it names a standard fact of the same chapter that the counted argument uses, and every step of the counted proof built on it is retained verbatim in the delivered body.\par\medskip
\ifdefined\hypertarget\hypertarget{sourcepointer-3}{}\fi
\noindent\textit{Source reference.} Original lemma 094L: sources/stacks-chow/tags.txt, source lines 11554--11554. Prerequisite pointer invoked at proof line(s) 12463; it names a standard fact of a sibling chapter that the counted argument uses, and every step of the counted proof built on it is retained verbatim in the delivered body.\par\medskip
\ifdefined\hypertarget\hypertarget{sourcepointer-4}{}\fi
\noindent\textit{Source reference.} Original lemma 093C: sources/stacks-chow/tags.txt, source lines 11510--11510. Prerequisite pointer invoked at proof line(s) 12440; it names a standard fact of a sibling chapter that the counted argument uses, and every step of the counted proof built on it is retained verbatim in the delivered body.\par\medskip
\ifdefined\hypertarget\hypertarget{sourcepointer-5}{}\fi
\noindent\textit{Source reference.} Original lemma 0901: sources/stacks-chow/tags.txt, source lines 11394--11394. Prerequisite pointer invoked at proof line(s) 12465; it names a standard fact of a sibling chapter that the counted argument uses, and every step of the counted proof built on it is retained verbatim in the delivered body.\par\medskip
\begin{lemma}
\label{lemma-category-constructible-modules}
Let $X$ be a quasi-compact and quasi-separated scheme. Let $\Lambda$
be a Noetherian ring. The category of constructible sheaves of
$\Lambda$-modules is exactly the category of modules of the form
$$
\Coker\left(
\bigoplus\nolimits_{j = 1, \ldots, m} j_{V_j!}\underline{\Lambda}_{V_j}
\longrightarrow
\bigoplus\nolimits_{i = 1, \ldots, n} j_{U_i!}\underline{\Lambda}_{U_i}
\right)
$$
with $V_j$ and $U_i$ quasi-compact and quasi-separated objects of
$X_\etale$. In fact, we can even assume $U_i$ and $V_j$ affine.
\end{lemma}

\begin{proof}
In the proof of Lemma \ref{lemma-torsion-colimit-constructible}
we have seen modules of this form are constructible. Since the
category of constructible modules is abelian
(Lemma \ref{lemma-constructible-abelian})
it suffices to prove that given a constructible module $\mathcal{F}$
there is a surjection
$$
\bigoplus\nolimits_{i = 1, \ldots, n} j_{U_i!}\underline{\Lambda}_{U_i}
\longrightarrow \mathcal{F}
$$
for some affine objects $U_i$ in $X_\etale$. By
Modules on Sites, Lemma
\ref{sites-modules-lemma-module-filtered-colimit-constructibles}
there is a surjection
$$
\Psi :
\bigoplus\nolimits_{i \in I} j_{U_i!}\underline{\Lambda}_{U_i}
\longrightarrow
\mathcal{F}
$$
with $U_i$ affine and the direct sum over a possibly infinite
index set $I$. For every finite subset $I' \subset I$ set
$$
T_{I'} = \text{Supp}(\Coker(
\bigoplus\nolimits_{i \in I'} j_{U_i!}\underline{\Lambda}_{U_i}
\longrightarrow \mathcal{F}))
$$
By the very definition of constructible sheaves, the set $T_{I'}$
is a constructible subset of $X$. We want to show that $T_{I'} = \emptyset$
for some $I'$. Since every stalk $\mathcal{F}_{\overline{x}}$ is
a finite type $\Lambda$-module and since $\Psi$ is surjective, for
every $x \in X$ there is an $I'$ such that $x \not \in T_{I'}$.
In other words we have
$\emptyset = \bigcap_{I' \subset I\text{ finite}} T_{I'}$. Since
$X$ is a spectral space by Properties, Lemma
\ref{properties-lemma-quasi-compact-quasi-separated-spectral}
the constructible topology on $X$ is quasi-compact by
Topology, Lemma \ref{topology-lemma-constructible-hausdorff-quasi-compact}.
Thus $T_{I'} = \emptyset$ for some $I' \subset I$ finite
as desired.
\end{proof}

\end{document}
